\documentclass[11pt]{article}
\usepackage[margin=1.05in]{geometry}
\usepackage{amsmath,amssymb}
\usepackage{booktabs}
\usepackage{longtable}
\usepackage{hyperref}
\hypersetup{colorlinks=true, linkcolor=blue!60!black, urlcolor=blue!60!black}

\title{The DRO\,$\rightarrow$\,Tulip Costate Libraries\\
\large Reference for \texttt{costate\_lib\_dro\_tulip.mat} (v1) and
\texttt{costate\_lib\_dro\_tulip\_v2.mat} (v2)}
\author{M.~Casey \and D.~Koblick}
\date{August 5, 2026}

\begin{document}
\maketitle

\begin{abstract}
Two MATLAB libraries of converged minimum-time low-thrust transfer solutions
between a distant retrograde orbit (DRO) and a tulip orbit in the Earth--Moon
circular restricted three-body problem. Each entry is a \emph{root} of the
indirect (Pontryagin) boundary-value problem --- initial costates and flight
time --- indexed by orbital phasing and, in v2, by thrust level. Every entry
is verified three ways, including acceptance by the production single-shooting
solver \texttt{pumpkyn.cr3bp.tfMin}, which returns library entries unchanged.
Version~1 holds 126 entries at a single operating point (0.07~N, $I_{sp}$
900~s); version~2 holds 990 entries spanning nine thrust levels from 1 to
15~N at $I_{sp}$ 1710~s.
\end{abstract}

\section{Purpose}

The indirect formulation of the minimum-time transfer reduces to finding eight
unknowns
\begin{equation}
z_8 \;=\; [\,\lambda_r(3);\;\lambda_v(3);\;\lambda_m;\;t_f\,]
\end{equation}
such that the trajectory propagated under the PMP control law satisfies
\begin{equation}
r(t_f) = r_f,\qquad v(t_f) = v_f,\qquad \lambda_m(t_f) = 0,\qquad H(t_f) = 0,
\label{eq:bcs}
\end{equation}
the last condition being free-final-time transversality. These are precisely
the equations \texttt{pumpkyn.cr3bp.tfMin} solves; the libraries store their
roots.

Single shooting on \eqref{eq:bcs} converges in seconds from a good guess and
fails outright otherwise --- the guess is the whole difficulty. A costate
library converts every future solve into a lookup plus a seconds-long polish,
and supplies continuation anchors for neighboring cases.

\section{Summary of the two libraries}

\begin{center}
\begin{tabular}{lll}
\toprule
 & \textbf{v1} & \textbf{v2} \\
\midrule
Structure name      & \texttt{costate\_lib\_dro\_tulip} & \texttt{costate\_lib\_dro\_tulip\_v2} \\
Axes                & departure $\times$ arrival phase & departure $\times$ arrival $\times$ \textbf{thrust} \\
Entries             & 126 & 990 \\
Phase-pair coverage & 126 of 144 & 130 of 144 (49 with all 9 rungs) \\
Thrust              & 0.07 N & 1, 1.5, 2, 3, 5, 7, 10, 12, 15 N \\
$I_{sp}$ / mass     & 900 s / 150 kg & 1710 s / 150 kg \\
Flight times        & 16.21--33.31 d & 0.269--2.560 d \\
Lookup helper       & \texttt{costate\_lib\_nearest.m} & \texttt{costate\_lib\_pick.m} \\
Worked example      & \texttt{costate\_lib\_example.m} & \texttt{costate\_lib\_v2\_example.m} \\
\bottomrule
\end{tabular}
\end{center}

\section{The two orbits}

Identical in both libraries and recorded inside each file, so a recipient
needs nothing but pumpkyn and the \texttt{.mat}.

\paragraph{Departure: DRO, selected by period.} \texttt{getDRO(tau)} returns
the distant retrograde orbit whose \emph{period} equals \texttt{tau} in
nondimensional time. Both libraries use $\tau = 1.0$, a period of
4.4327~days.

\paragraph{Arrival: tulip.} \texttt{getTulip(tau, Np, pm)} with
$\tau = 5\cdot 2\pi/6$, $N_p = 7$, $pm = -1$; period 23.2093~days. The tulip
period is exactly $5/6$ of the Moon's revolution, so the seven-lobed pattern
closes after six tulip periods.

\paragraph{Phase convention.} A phase is the fraction of an orbit period past
the reference state returned by \texttt{getDRO}/\texttt{getTulip}; entries
store it as a fraction, in nondimensional time, and in days. Both axes wrap,
so the domain is a torus. The arrival axis is offset by the anchor phase
$0.075378$ in both libraries, so their cells correspond one-to-one. The grid
is $12 \times 12 = 144$ phase pairs, spaced 8.9~hours along the DRO and
1.93~days along the tulip.

\section{Constants and thruster conversions}

\begin{align}
\mu^\ast &= 0.012150585609624, &
l^\ast &= 389703.264829278~\text{km}, \\
t^\ast &= 382981.289129055~\text{s}, &
1~\text{ND time} &= 4.432654~\text{days}.
\end{align}

\texttt{tfMin} consumes nondimensional thrust and exhaust velocity:
\begin{equation}
T_{\max}^{\text{nd}} = \frac{T[\mathrm{N}]}{m_0[\mathrm{kg}]}\,
\frac{(t^\ast)^2}{l^\ast \cdot 1000},
\qquad
c^{\text{nd}} = \frac{I_{sp}}{t^\ast}\,g_0,
\qquad
g_0 = \frac{9.80665\,(t^\ast)^2}{1000\,l^\ast}.
\end{equation}
Mass is normalized so $m(0)=1$; multiply by \texttt{m0\_kg} for kilograms.

Version~1 is a single operating point (0.467~mm/s$^2$ initial acceleration,
transfers spanning several DRO revolutions). Version~2 spans
6.7--100~mm/s$^2$, where every transfer is \emph{sub-revolution}. The two are
neighboring sheets rather than points on one continuum: the interval between
0.07~N and 1~N contains several revolution-count transitions, which is why v2
was anchored at high thrust instead of continued upward from v1.

\section{Structure layout}

Both files contain a single struct with fields \texttt{name},
\texttt{description}, \texttt{created}, \texttt{provenance},
\texttt{constants}, \texttt{thruster}, \texttt{departure\_orbit},
\texttt{departure\_params}, \texttt{arrival\_orbit}, \texttt{arrival\_params},
\texttt{entries}, \texttt{n\_entries}, \texttt{grid}, \texttt{usage}.

\begin{center}
\begin{tabular}{ll}
\toprule
\textbf{\texttt{entries(k)} field} & \textbf{Meaning} \\
\midrule
\texttt{departure\_phase\_frac/\_nd/\_days} & where on the DRO the transfer starts \\
\texttt{arrival\_phase\_frac/\_nd/\_days}   & where on the tulip it ends \\
\texttt{thrust\_N}, \texttt{isp\_s}, \texttt{m0\_kg} & the propulsion case for this entry \\
\texttt{lambda0} $[7\times1]$ & initial costates $[\lambda_r;\lambda_v;\lambda_m]$ \\
\texttt{tf\_nd}, \texttt{tf\_days} & minimum flight time \\
\texttt{z8} $[8\times1]$ & $[\lambda_0; t_f]$, the \texttt{tfMin}-ready vector \\
\texttt{ms\_residual} & multiple-shooting residual at convergence \\
\texttt{flight\_miss\_km} & arrival error when the PMP law is flown end to end \\
\bottomrule
\end{tabular}
\end{center}

The \texttt{grid} substructure repeats the same information as lookup arrays:
phase axes, \texttt{thrust\_N} (v2), a \texttt{has\_solution} mask,
\texttt{tf\_days}, and \texttt{entry\_index} into \texttt{entries}
($0 =$ no solution). Arrays are $12\times12$ in v1 and $12\times12\times9$ in
v2.

\section{Production pipeline}

Both libraries follow the same three steps (recorded in
\texttt{process/COSTATE\_LIBRARY\_PIPELINE.md}):

\begin{enumerate}
\item \textbf{Direct solve.} Hermite--Simpson collocation yields a trajectory
and a costate trajectory (Hager covector mapping of the NLP defect
multipliers), verified by flying the solved control end to end.
\item \textbf{Refinement.} \texttt{ms\_tfmin} --- multiple shooting with an
analytic block Jacobian assembled from pumpkyn's own segment state-transition
matrices --- drives the residual of \eqref{eq:bcs} to $\sim10^{-12}$.
Collocation duals are only $\sim10^{-3}$ accurate and single shooting
amplifies that by $\sim10^3$ over a multi-revolution arc, so this step is
mandatory.
\item \textbf{Acceptance.} \texttt{pumpkyn.cr3bp.tfMin} must return the entry
essentially unchanged.
\end{enumerate}

The two differ in grid strategy. Version~1 swept the torus in two waves ---
cold straight-line seeds opened isolated cells, then solutions flooded to
neighbors --- because a cold solve at 0.07~N is expensive and often fails.
Version~2 anchored each phase pair at 15~N, where the transfer is nearly
impulsive and converges cold in seconds, then walked thrust \emph{down} with
each rung warm-started from the one above. Continuation keeps a cell on a
single solution family; independent cold solves do not (two meshes at 15~N
were measured landing 50\,\% apart in $t_f$, both feasible).

\section{Verification}

\begin{center}
\begin{tabular}{lll}
\toprule
\textbf{Check} & \textbf{v1 (126 entries)} & \textbf{v2 (990 entries)} \\
\midrule
Multiple-shooting residual & max $9.5\times10^{-11}$ & max $9.3\times10^{-11}$ \\
Flown arrival miss & max $0.064$ km & max $1.5\times10^{-5}$ km \\
\texttt{tfMin} acceptance $\lVert\Delta z\rVert$ & max $2.8\times10^{-8}$, 126/126 & max $1.2\times10^{-10}$, 990/990 \\
\bottomrule
\end{tabular}
\end{center}

For v1 there is an additional cross-method check: direct and indirect flight
times agree to a median of $2\times10^{-8}$ ND over all 126 phase pairs.

\section{Coverage and known gaps}

Version~1 solves 126 of 144 phase pairs; the missing 18 are dominated by the
arrival row at phase $0.0754$, the tulip's fastest-moving point, which is hard
to reach after four to seven revolutions.

Version~2 covers 130 of 144 pairs, 49 of them with the complete ladder:

\begin{center}
\begin{tabular}{rrl}
\toprule
\textbf{Thrust (N)} & \textbf{Pairs} & \textbf{$t_f$ range (days)} \\
\midrule
15   & 109 & 0.269--0.661 \\
12   & 128 & 0.301--0.583 \\
10   & 128 & 0.330--0.642 \\
7    & 121 & 0.398--0.773 \\
5    & 125 & 0.483--0.919 \\
3    & 119 & 0.704--1.375 \\
2    & 106 & 0.923--1.553 \\
1.5  &  91 & 1.122--2.054 \\
1    &  63 & 1.546--2.560 \\
\bottomrule
\end{tabular}
\end{center}

Two patterns stand out. The 15~N anchor is cold-started and therefore fails
more often (109 pairs) than the rungs beneath it (128); a failed anchor does
not stop the ladder, since the next rung recovers. The bottom rung is thin ---
at 1~N only 63 pairs converge --- and filling it is the main outstanding item,
planned as an intermediate 1.25~N rung to soften the continuation step.
Version~2 does, however, cover phase pairs version~1 could not: the hard
arrival row is reachable at sub-revolution thrust, confirming that its
difficulty was many-revolution geometry rather than the endpoints themselves.

\section{Usage}

\begin{verbatim}
L = load('costate_lib_dro_tulip_v2.mat');  lib = L.costate_lib_dro_tulip_v2;
[tf_days, e, bracket] = costate_lib_pick(lib, depDays, arrDays, thrustN);

% e.z8 is a converged seed at the nearest rung; tf_days is interpolated
% linearly between the bracketing rungs. Rebuild both orbits from
% lib.departure_params / lib.arrival_params, then either fly it:
[t, y] = pumpkyn.cr3bp.tfMinProp(e.tf_nd, [rv0(1:6)'; 1; e.lambda0], ...
                                 Tnd, lib.thruster.c_nd, mu);
% ...or hand the seed to the single-shooting solver, which returns it
% unchanged:
z = pumpkyn.cr3bp.tfMin(rv0(1:6), rvf(1:6), e.z8, Tnd, lib.thruster.c_nd, mu);
\end{verbatim}

Version~1 is identical with \texttt{costate\_lib\_nearest(lib, depDays,
arrDays)}, which takes no thrust argument. For a thrust or phase between grid
values, the returned \texttt{z8} is a warm start a fraction of a step away and
\texttt{tfMin} closes the difference in about a second. To browse rather than
query, \texttt{lib.grid.has\_solution} gives availability,
\texttt{lib.grid.tf\_days} the flight times, and
\texttt{lib.grid.entry\_index} maps a grid cell to its row in
\texttt{lib.entries}.

\section{The multi-orbit catalog (August 6)}

The pipeline swept the admissible box: 16 sheets --- DRO periods
$\tau \in \{0.5, 1, 2, 3\}$ ND $\times$ petal counts $N_p \in \{5,7,9,12\}$
($pm=-1$; the $+1$ branch is the exact $z$-mirror) --- at $6\times6$ phasing
and nine thrust rungs, yielding \textbf{3{,}936 three-gate entries} with
89--97\,\% phase-pair coverage on every sheet (\texttt{costate\_catalog\_%
dro\_tulip.mat}, compact ND-only format). Cross-sheet findings: at fixed
5~N, the orbit-pair choice spans $\Delta V$ 1.18--4.09~km/s ($3.5\times$,
larger than phasing's $\sim2\times$); the fastest transfer anywhere is
0.218~d ($\tau=0.5$, $N_p=9$, 15~N) and the slowest 3.44~d ($\tau=3$,
$N_p=12$, 1~N) at nearly equal $\Delta V$.

\paragraph{Halo bounds (the next family).} Applying the same admissibility
survey to pumpkynPie's Halo family: \emph{L2 southern} is continuously
admissible for periods 7.1--15.1 days (periselene 2{,}800--49{,}900~km);
L1 is largely sub-surface (two clean members); L3 lies 0.39--1.05~Gm from
the Moon, outside the 100~Mm box. The Halo$\to$tulip campaign proceeds on
L2 southern via the shared \texttt{get\_family\_orbit} provider.

\section{Growth axes}

The pipeline is parameterized, so further libraries are reruns rather than
rebuilds: \emph{DRO period} (the selector $\tau$ \emph{is} the period, so each
$\tau$ yields one sheet, and together they form the $\Delta V$--period--phasing
map), \emph{thrust} (further rungs or ranges, as v2 demonstrates), and
\emph{orbit pair} (endpoints from any of pumpkyn's catalogued families).

\end{document}
